\section{Measuring Premise Responses}
\label{sec:method}

\subsection{Problem setup}
A question such as \textit{``What has $x$ done since becoming the CEO of $y$?''} presupposes $f=\langle x,\textit{role},y\rangle$: person $x$ holds a role at entity $y$. Let $t_f$ denote the date when this became true. We ask how a model responds when $f$ is true at the benchmark snapshot but $t_f$ falls after its training data. We label a response \textbf{accept} if it answers as though $f$ holds, \textbf{rebut} if it denies $f$ or names someone else as the current holder, and \textbf{cannot verify} if it does neither. We also call the last outcome a \emph{deferral}.

To distinguish reduced rebuttals from improved discrimination, we compare the newly-true holder (\NEW{}) with a former holder (\ANC{}) and a holder swapped from another entity (\SWP{}):
\begin{align*}
J_{\text{fam}} &= P(\text{rebut}\mid\ANC{})-P(\text{rebut}\mid\NEW{}),\\
J_{\text{unf}} &= P(\text{rebut}\mid\SWP{})-P(\text{rebut}\mid\NEW{}).
\end{align*}
These margins, reported in percentage points, are positive when false premises are rebutted more often. Zero means no separation, whether both rates are high or low. The subscripts denote intended familiar/unfamiliar controls, not verified familiarity (\S\ref{sec:analysis}).

\subsection{Items from Wikidata}
We query Wikidata \citep{vrandecic2014wikidata} for statements with start (P580) and end (P582) qualifiers on five relations: head of government (P6), head of state (P35), CEO (P169), head coach (P286) and chairperson (P488), and freeze the dump at 2026-11-10. A \emph{transition} is an entity whose single current holder replaced a predecessor whose term ended within 400 days. We group the items by the timing and truth of their premises. \NEW{} (true): the current holder of a transition; \NEW{}-post items became true on or after 2025-01-01, after the public release of every model we pool (\S\ref{sec:setup}), and \NEW{}-near items changed between mid-2023 and 2024. \OLD{} (true): a holder in office since 2021 or earlier with no recorded change since 2022. \STRE{} and \STRL{} (false): a holder replaced before 2021, or in 2021--mid-2022. \STP{} (false): the predecessor of a \NEW{}-post holder; \NEV{} (false): the long-standing holder of a different entity with the same relation, often implausible on its face. Each of these items is asked with one of three templates, assigned in rotation before the audit, so the counts per template are not exactly balanced (Appendix~\ref{app:prompts}). Two of the templates are historically true of a former holder, so the false groups of this grid mix ``never held the role'' with ``no longer holds it''; we therefore also report them on the present-tense template alone (\S\ref{sec:results}).

\paragraph{Entity triples (main test).}
Our main test uses the present-tense template \textit{``What are the main priorities of $x$ as the current $r$?''}. Each triple substitutes three people: \NEW{}, the newly-true holder; \ANC{}, the most recent former holder of the same role at the same entity recorded in our Wikidata dump whose last term ended before 2016 and who is alive (no recorded death); and \SWP{}, the \NEW{}-post holder of another entity with the same relation (a within-relation derangement, assigning every holder to a different entity). Like \NEW{}, the swapped person took a role recently, but this premise is false. Popularity is the number of Wikidata sitelinks (linked Wikimedia pages). Median counts are \WSlS{} for \SWP{}, \WSlN{} for \NEW{} and \WSlA{} for \ANC{}, but the match holds only in the median: in \WSwMoreFam{} of the \WNTri{} triples the swapped person is more popular than the new holder, and most swaps cross countries (\S\ref{sec:analysis}). Entity, role phrase and wording are identical within a triple. \WNTri{} complete triples remain after the audit.

\paragraph{Validation and audit.}
We require Wikipedia corroboration on the snapshot date: the current holder's surname in the entity article's lead or infobox. Office labels must match a ``position held'' (P39) statement. For the subsequent audit, Claude Haiku 5.5 flagged truth labels that contradicted its knowledge; we checked these against each person's Wikipedia article. We also checked every \NEW{} holder's lead for acting or interim roles, removed reinstated holders whom models might already know in the role (e.g.\ Donald Tusk), and checked reviewer flags (e.g.\ TASS's head, appointed in 2023, not 2012). We removed \VNExcl{} items and release the audit log (Appendix~\ref{app:data}).

\subsection{Interventions}
\label{sec:interventions}
All conditions use identical items (prompts in Appendix~\ref{app:prompts}). \textsc{Plain} asks the question alone. \textsc{Premise-check} adds our instruction to check the question for a false premise and correct it; \textsc{PC (FreshQA)} adds FreshQA's exact sentence, ``Please check if the question contains a valid premise before answering'' \citep{vu2024freshllms}. \textsc{Date} gives today's date; \textsc{Cutoff-stated} adds that the model's training data ends before 2025 (true for every pooled model). \textsc{Policy} is our cutoff-aware instruction plus the date: accept premises consistent with the model's knowledge or explainable by a post-cutoff change, rebut only contradictions such a change cannot explain, and otherwise say the premise cannot be verified; \textsc{Policy-nodate} drops the date. \textsc{RAG-fresh} prepends an oracle snippet stating the current holder and start date, and \textsc{RAG-fresh+PC} adds our premise-check instruction to it, as in FreshQA's setting. \textsc{RAG-stale} renders the snippet from the state before the most recent change, and \textsc{RAG-mixed} gives the stale snippet, labelled as an older cached snapshot, and the dated fresh one in alternating order.

\subsection{Cutoff and Belief Probes}
\label{sec:cutoffs}
Reported cutoffs are often missing or inaccurate \citep{cheng2024dated}. For each model, we ask \textit{``Who is the $r$? Answer with the person's name only.''} for 439 transitions spread over 2019--2026, record whether the new holder's name is more probable than the old one's, and fit a step function by maximum likelihood (bootstrap confidence interval, CI; permutation test). We use these estimates only descriptively: they identify which models know office holders at all, and they are sensitive to how name probabilities are normalised (Appendix~\ref{app:extra}). Our \NEW{}-post items do not rely on them, because they postdate every pooled model's release. We also score a yes/no belief probe, the probability $P(\textit{No})$ assigned to a negative answer to \textit{``Is it true that $x$ is the $r$?''}.

\subsection{Response labelling}
\label{sec:labels}
Three large language models (LLMs) label every response with the same prompt: Claude Haiku 5.5 (primary; an API model that is not evaluated), Qwen2.5-7B-Instruct and Mistral-7B-Instruct-v0.3 (both also evaluated). They choose among accept, rebut, cannot verify and refuse; refusals are merged into cannot verify, and we report them separately where it matters. \emph{Rebut} includes hedged denials (``as of my last update, $x$ is not the CEO''); we report hedged and unhedged rebuttals separately (\S\ref{sec:analysis}). Haiku was chosen as primary judge on a first validation sample (\VSilN{} responses labelled by Claude Sonnet 5.5; Appendix~\ref{app:judge}). For held-out validation, Claude Sonnet 5.5 (the primary judge's family) and OpenAI gpt-5.3-codex (another family) label all seven models' triple responses under \textsc{Plain}, \textsc{Premise-check}, \textsc{Cutoff-stated} and \textsc{RAG-fresh+PC}. We report agreement and repeat the main test under each labeller (\S\ref{sec:analysis}). No human labels were collected; Appendix~\ref{app:human} gives the pending protocol.
